\section{From geographic sales to legal entities}
\label{sec:apple-perimeter}

A regional sales forecast still leaves open which company sells the goods,
owns the inventory or owes the supplier. Those identities matter when a
scenario changes a particular entity's tax, funding or supply arrangements.
Apple's geographic segments cannot answer that question: their sales are
generally attributed by customer location and retail-store geography. The
consolidation policy and subsidiary exhibit supply a different kind of
information, about the companies represented in the group statements.

Annual Note 1 says that the consolidated financial statements include the
accounts of Apple Inc. and its wholly owned subsidiaries. The filing also
defines ``Company'' and ``Apple'' collectively in that way, unless otherwise
stated. Interim Note 1 repeats the wholly-owned scope for its condensed
statements and requires them to be read with the fiscal 2024 annual statements
and notes \citep[definition of Company and Note 1, printed pp.~1 and 34]{apple2024}
\citep[Note 1, printed p.~5]{apple2025q1}. This is positive evidence about the
reported consolidation perimeter. It is not a set of separate parent-company
accounts, nor does it identify every legal entity or ownership link within
that perimeter.

Exhibit 21.1 names the following subsidiaries and their jurisdictions of
incorporation. The labels below are transcribed from the exhibit, including
the separate U.S. state jurisdictions \citep[Exhibit 21.1, retained PDF p.~114]{apple2024}.
\begin{center}
\small
\begin{tabular}{ll}
\toprule
Subsidiary & Jurisdiction of incorporation\\\midrule
Apple Asia Limited & Hong Kong\\
Apple Asia LLC & Delaware, U.S.\\
Apple Canada Inc. & Canada\\
Apple Computer Trading (Shanghai) Co., Ltd. & China\\
Apple Distribution International Limited & Ireland\\
Apple India Private Limited & India\\
Apple Insurance Company, Inc. & Arizona, U.S.\\
Apple Japan, Inc. & Japan\\
Apple Korea Limited & South Korea\\
Apple Operations International Limited & Ireland\\
Apple Operations Limited & Ireland\\
Apple Operations Mexico, S.A. de C.V. & Mexico\\
Apple Pty Limited & Australia\\
Apple Sales International Limited & Ireland\\
Apple South Asia (Thailand) Limited & Thailand\\
Apple Vietnam Limited Liability Company & Vietnam\\
Braeburn Capital, Inc. & Nevada, U.S.\\
iTunes K.K. & Japan\\\bottomrule
\end{tabular}
\end{center}

The exhibit expressly limits the list. Its footnote states: ``Pursuant to
Item 601(b)(21)(ii) of Regulation S-K, the names of other subsidiaries of Apple
Inc. are omitted because, considered in the aggregate, they would not
constitute a significant subsidiary as of the end of the year covered by this
report.'' This preserves Apple's stated reason and date for the omission;
it does not establish a numerical significance test for this model. The
eighteen rows cannot be treated as an exhaustive group roster or as proof that
an omitted subsidiary is irrelevant to every entity-specific question.

The same filing makes that incompleteness concrete. Its tax-risk discussion
mentions Ireland and Singapore, ``where a number of the Company's subsidiaries
are organized,'' yet no Exhibit 21.1 row has Singapore as its jurisdiction
\citep[Risk Factors, printed p.~16]{apple2024}. The evidence supports
subsidiaries organized there while leaving their names and ownership paths
unresolved. Creating a company named ``Apple Singapore'' would turn a
geographic reference into an invented identity.

\subsection{Membership does not determine the ownership path}

Wholly-owned membership and immediate ownership answer different questions.
For a simple hypothetical example, a subsidiary could be owned directly by
the parent, or through another subsidiary that the parent wholly owns. If
every link is wholly owned, both arrangements preserve ultimate wholly-owned
membership while producing different immediate parents. This logical example
does not select either arrangement for any company in the table. Neither the
similar names of Apple Operations Limited and Apple Operations International
Limited nor their common incorporation jurisdiction establishes an ownership
link between them.

A model can therefore retain each named subsidiary, its reported jurisdiction
and the general consolidation assertion without drawing unsupported edges.
It should record an immediate parent, share class, ownership interest and
effective interval only when the evidence identifies them. The September 2024
exhibit is not a control history for earlier years or an updated list at
28 December 2024. Repetition of the general policy in the interim note does
not show that every subsidiary and ownership path remained unchanged.
Determining the dated group requires incorporation, acquisition, disposal,
reorganisation and control evidence as relevant; it cannot be recovered from
the fiscal calendar alone.

\subsection{Selling, manufacturing and owning are separate relationships}

Apple India Private Limited illustrates the distinction between incorporation
and sales attribution. The exhibit places its incorporation in India, while
Apple's segment definition places India in Europe. Under the general
customer-location rule, an India-customer sale belongs in that sales region;
neither statement identifies its legal seller or assigns all India sales to
that subsidiary. Even a complete subsidiary list would not supply the missing
seller-by-customer-region relationship.

The Business discussion also separates distribution channels: direct sales
through Apple's stores and sales force accounted for 38\% of fiscal 2024 net
sales, and indirect channels for 62\%. The latter include third-party cellular
carriers, wholesalers, retailers and resellers
\citep[Markets and Distribution, printed p.~2]{apple2024}. These are sales
channel shares, not ownership shares or counts of companies. A trading
relationship does not by itself establish that a counterparty belongs in the
consolidated group.

Production adds a further distinction. Apple says substantially all its
manufacturing is performed in whole or in part by outsourcing partners,
primarily in China mainland, India, Japan, South Korea, Taiwan and Vietnam.
It has also invested in manufacturing process equipment, much of which is
held at certain outsourcing partners, and made prepayments under long-term
supply agreements \citep[Risk Factors, printed p.~8]{apple2024}. Its MD\&A
describes partners manufacturing subassemblies and performing final assembly
and testing, alongside purchases from component suppliers
\citep[Manufacturing Purchase Obligations, printed p.~25]{apple2024}.
Equipment custody, a supply contract and corporate control are thus separate
relationships. A disruption at a partner may affect product supply, equipment
recoverability and prepayments without supplying evidence that the partner
should be consolidated. The affected assets, contracts and counterparties
must be identified before amounts can be assigned to that scenario.

\subsection{Tax and currency queries need facts beyond incorporation}

The jurisdiction column also leaves entity tax residence and functional
currency unresolved. Functional currency is the currency in which an entity
measures its accounts; it is a different question from the dollars in which
Apple presents its group statements or the currency of a particular invoice.
The group foreign-exchange discussion describes Apple as a net receiver of
currencies other than the U.S. dollar, with pricing and hedging responses to
exchange-rate movements. It supplies no entity-by-entity functional-currency
assignment \citep[Foreign Exchange Rate Risk, printed p.~27]{apple2024}.
An Ireland incorporation row therefore does not justify entering euros as the
entity's functional currency, just as the India row does not establish rupees.

Tax evidence has its own dates and units of analysis. Annual Note 7 describes
the historic tax allocation of profits of Irish branches of two subsidiaries
without naming those subsidiaries in that passage. The MD\&A describes the
resulting group payment obligation and restricted escrow
\citep[Note 7, printed pp.~39--40; State Aid Decision Tax Payable, printed p.~26]{apple2024}.
The four Ireland rows in Exhibit 21.1 do not identify which entities had those
branches, their historic tax residence or the jurisdictional tax bases used
in current accounts. A name match or country match would bypass the legal
and temporal facts needed for the tax calculation.

These distinctions determine what a supported query can return. The retained
sources answer questions about the stated consolidated scope and the named
subsidiaries' incorporation jurisdictions. Entity-specific tax, currency and
intercompany questions require additional records: legal counterparties,
dated ownership and control, compatible entity books, functional currencies,
tax bases and matched transactions. Group eliminations then require the
applicable accounting methods as well as those facts. The later IFRS group
examples explain their own assumptions; they do not establish the operative
US GAAP treatment for Apple. Keeping those dependencies explicit preserves
useful source knowledge while leaving the unanswered entity questions open.
